\documentclass[12pt]{article}
\usepackage[utf8]{inputenc}
\usepackage[T1]{fontenc}
\usepackage{hyperref}
\title{Tapiafaru Language Resources for K-12 Curriculum Development}
\author{Compiled by Claude Sonnet 5.5 (claude-sonnet-5-5) for Andr\'as Kornai}
\date{}
\begin{document}
\maketitle

\section{Introduction}
Tapiafaru (ISO 639-3: ptr; Glottocode piam1242), also called Piamatsina, is an Oceanic (Austronesian) language spoken in the north of Espiritu Santo, Vanuatu. Estimates of speakers are small (150 in an older 1981 figure, 250 in another source; Joshua Project gives a population of 400). It should not be confused with Futuna-Aniwa, a separate Polynesian language of Tafea Province. Very little material specific to this language was found.

\section{Primary Educational Resources}
No primers, textbooks or school programmes in Tapiafaru were found in the searches.

\section{Linguistic and Reference Materials}
\begin{itemize}
\item Wikipedia article on the Piamatsina (Tapiafaru) language: \url{https://en.wikipedia.org/wiki/Piamatsina_language}
\item ``The languages of Vanuatu'' (ANU Open Research repository), a general survey of Vanuatu languages: \url{https://openresearch-repository.anu.edu.au/bitstream/1885/14819/5/The%20languages%20of%20Vanuatu%202015.pdf}
\end{itemize}

\section{Religious and Community Resources}
\begin{itemize}
\item Joshua Project language page for ptr (scripture status listed as translation needed): \url{https://joshuaproject.net/languages/ptr}
\end{itemize}

\section{Digital and Audio Resources}
No audio, video, corpus or software resources specific to Tapiafaru were found in the searches.

\section{Curriculum Development Context}
The language is spoken by the whole small community but is not known to be taught in schools. Only overview-level sources were found, and no dictionary, grammar or scripture is evident. Curriculum work would need documentation to be undertaken first, together with the community.

\end{document}